% Standalone preview. Copy only the Proposed Method section into the manuscript.
% Citation keys match docs/references.bib; the preview bibliography is embedded.
% Notation aligned with the author's manuscript supplied on 2026-10-06.
\documentclass[10pt,twocolumn]{article}
\usepackage[a4paper,margin=19mm]{geometry}
\usepackage{amsmath,amssymb}
\usepackage[hidelinks]{hyperref}
\begin{document}
\setcounter{section}{3}

\section{Proposed Method}
\label{sec:method}

\subsection{Decision Process and Hierarchical Framework}
\label{subsec:framework}

Server placement and frequency control are coordinated through a
hierarchical partially observable decision process. The two learners
observe task and resource features rather than the complete scheduling
state. We describe the process by
\begin{equation}
 \mathfrak P=\langle\mathcal K,\mathcal S,\mathcal A,\mathcal P,
 \Omega,\mathcal O,\mathbf r,\boldsymbol\delta\rangle,
 \label{eq:method_pomdp}
\end{equation}
where:
\begin{itemize}
 \item $\mathcal K=\{H,L\}$ is the set of upper and lower learners.
 \item $\mathcal S$ is the global state space, including workflow progress,
 server queues, current time, and the arrival-process state.
 \item $\mathcal A$ is the joint action space. For a ready task $i$,
 $a_t=(m_i,\rho_i)$ specifies its server and frequency ratio, with
 $\mathcal A=\bigcup_{m\in\mathcal M}
 (\{m\}\times[\underline\rho_m,1])$.
 \item $\mathcal P(s'\mid s,a)$ is the state-transition kernel induced by
 task execution, queue evolution, and workflow arrivals.
 \item $\Omega=\Omega_H\times\Omega_L$ is the joint observation space,
 where $\Omega_H$ and $\Omega_L$ contain the feature representations
 $s_t^H$ and $s_t^L$, respectively.
 \item $\mathcal O$ is the observation kernel. Upper observations are
 obtained from the current state; lower observations additionally depend
 on the server selected by the upper learner. The encoders used here
 are deterministic.
 \item $\mathbf r=(r_{E,t}^H,r_t^L)$ contains the upper energy reward and
 lower frequency-control reward. Deadline and reliability costs are
 specified separately in the following subsections.
 \item $\boldsymbol\delta=(\delta_E,\delta_D,\delta_R,\delta_L)$ contains
 the discount factors for upper energy and constraint estimation and
 lower actor--critic learning, each taking a value in $(0,1]$.
\end{itemize}

The index $t$ denotes a task-dispatch decision, not a fixed physical-time
interval. A common ready-task ordering rule supplies task $i$ of workflow
$w$; task ordering is not learned. The upper learner first selects $m_i$
from $s_t^H$, after which the lower learner determines $\rho_i$ from the
selected-server observation $s_t^L$. Each trajectory covers the complete
execution of a workload.

Placement determines communication and contention, whereas frequency
control determines execution time and active energy on the selected
server. We therefore factorize the policy into a discrete upper policy
and a conditional continuous lower policy:
\begin{equation}
 \pi(m_i,\rho_i\mid s_t)
 =\pi_H(m_i\mid s_t^H)\,
   \pi_L(\rho_i\mid s_t^L,m_i),
 \label{eq:method_factorization}
\end{equation}
where $s_t^H$ and $s_t^L$ are the respective feature-based observations.
They summarize, rather than fully reconstruct, the Markov state.
The upper learner uses three D3QN estimators for energy, deadline, and
reliability, while the lower learner uses soft actor--critic (SAC) for
DVFS. Both learn from the same trajectories: the upper action is evaluated
after the lower frequency decision, and both observe global deadline
pressure. Their constraint costs distinguish placement consequences from
the additional losses caused by frequency reduction.

For the current workflow $w$, retain the local task notation
$\mathcal N$ and $n=|\mathcal N|$ from the system model. Let
$\mathcal I_w$ contain its decision indices and $[x]_+=\max(0,x)$.
The service threshold $\Gamma$ and target probability $\eta$ retain
their meanings in the problem formulation.
In sums over decisions from different workflows, $n$ and $D_w$ refer
to the workflow of the corresponding decision.
We use $\mu$, rather than $\lambda$, for constraint multipliers to
distinguish them from the transient-fault rates in the system model.
The actual energy remains $E_i$; $\widehat E(i,m,\rho)$ denotes the same
energy model evaluated for a candidate server and frequency. The
energy-efficient ratio $\rho_m^{\mathrm{ee}}$ is defined in the power and
energy model.

\subsection{Upper-Level Server Selection}
\label{subsec:high}

\subsubsection{State Representation}

The upper state contains task workload and deadline slack, together with
server-specific estimates of communication time, execution time, queue
waiting, full-frequency energy, and task failure probability. Server
occupancy and recent assignment history describe resource contention.
Current slack alone does not account for unfinished successors. Let
$\ell_i$ be an estimate of the completion time still required after task
$i$, obtained from its upward rank after removing its mean execution
time; any required final output transfer is included. For candidate
server $m$, let $\widehat{ST}(i,m)$ and $\widehat{FT}(i,m)$ be its
estimated start and full-frequency finish times under the current queues.
Its predicted workflow response time is
\begin{equation}
 \widehat T(i,m)=\widehat{FT}(i,m)-a_w+\ell_i.
 \label{eq:method_predicted_finish}
\end{equation}
With a small margin $\zeta$, the available execution budget is
$b(i,m)=(1-\zeta)[a_w+D_w-\widehat{ST}(i,m)-\ell_i]$.
Let $\widehat{ET}(i,m)=L_i/F_m$ denote full-frequency execution time on candidate
server $m$. The state includes this budget and the indicative frequency
\begin{equation}
 \widehat\rho(i,m)=
 \begin{cases}
 \min\{1,\max\{\rho_m^{\mathrm{ee}},\widehat{ET}(i,m)/b(i,m)\}\},&b(i,m)>0,\\
 1,&b(i,m)\leq0,
 \end{cases}
 \label{eq:method_time_frequency}
\end{equation}
This frequency is a state feature, not a hard deadline-feasibility
guarantee.

To expose congestion beyond the current workflow, let $\mathcal W_t$
contain all released, unfinished workflows and let $h_w$ estimate their
remaining critical-path duration, accounting for the residual execution
of running tasks. At current physical time $t_{\mathrm{now}}$, and with
$\epsilon>0$ to prevent division by zero, define
\begin{align}
 p_w&=\frac{h_w}{\max\{a_w+D_w-t_{\mathrm{now}},\epsilon\}},\nonumber\\
 \chi_t&=\tanh\!\left(\frac{\sum_{w\in\mathcal W_t}p_w}{|\mathcal M|}\right),\quad
 \chi_t^{\max}=\frac{\min\{2,\max_{w\in\mathcal W_t}p_w\}}{2},
 \label{eq:method_global_pressure}
\end{align}
Both features are zero when $\mathcal W_t$ is empty. They enter both levels,
allowing placement and frequency control to observe shared deadline
pressure rather than only local slack. Time features are normalized by
the corresponding workflow deadline.

\subsubsection{Action Selection}

The upper learner maintains independent energy-return, deadline-cost,
and reliability-cost estimators, denoted by $Q_E$, $Q_D$, and $Q_R$.
Each uses the dueling representation~\cite{wang2016dueling}
\begin{equation}
 Q_h(s,m)=V_h(s)+A_h(s,m)
 -\frac{1}{|\mathcal M|}\sum_{m'\in\mathcal M}A_h(s,m'),
 \quad h\in\{E,D,R\}.
 \label{eq:method_dueling}
\end{equation}
The estimators have separate parameters; their outputs are combined
only for action selection. Given multipliers $\mu_D^H$ and $\mu_R^H$,
the exploitation action is
\begin{equation}
 m_i=\arg\max_{m\in\mathcal M}
 \left\{Q_E(s_t^H,m)-\mu_D^H[Q_D(s_t^H,m)]_+
 -\mu_R^H Q_R(s_t^H,m)\right\}.
 \label{eq:method_high_action}
\end{equation}
The nonnegative deadline projection prevents a negative cost estimate
from being treated as a scheduling benefit. Training uses
$\varepsilon$-greedy exploration over the available servers.

\subsubsection{Reward and Constraint Costs}

\paragraph{Energy.}
Let $\bar E_i=\operatorname{median}_{m\in\mathcal M}
\widehat E(i,m,1)$ be its full-frequency reference scale. The energy feedback is
\begin{equation}
 r_{E,t}^H=\frac{\bar E_i-E_i}{n\bar E_i}.
 \label{eq:method_high_energy}
\end{equation}
This dimensionless surrogate exposes the realized energy difference
while limiting the dependence of reward magnitude on workflow size.

\paragraph{Deadline.}
Let $\Delta T_{w,t}$ be the nonnegative increment in the predicted
partial workflow response time after decision $t$. Terminal feedback
reconciles these increments with the completed response time, such that
$\sum_{t\in\mathcal I_w}\Delta T_{w,t}=T_w$.
Let $q_t$ be the selected task's waiting time caused by server contention,
excluding prerequisite communication and waiting before dispatch.
It is the server-contention component of $WT(i)$, not the
communication delay $c_{ji}$ or the execution duration $ET_i$.
The dense cost is
\begin{equation}
 d_t=\frac{\Delta T_{w,t}}{D_w}
 +\beta_q\frac{q_t}{nD_w}.
 \label{eq:method_dense_cost}
\end{equation}
The queue term provides credit to tasks whose waiting does not increase
the current workflow completion frontier. It also participates in
distributing the final miss cost. Define
$b_w=\mathbb{I}(T_w>D_w)$ and
$\omega_t=d_t/\sum_{j\in\mathcal I_w}d_j$.
For a completed trajectory $e$ with workflow set $\mathcal W_e$, let
$\mathrm{SR}_e$ be its deadline success rate as defined in the problem
formulation. The service-level failure indicator is
$z_e=\mathbb{I}(\mathrm{SR}_e<\Gamma)$.
To allocate this trajectory-level failure, define the avoidable
predicted completion regret
$\Delta_t^{\mathrm{sel}}=[\widehat T(i,m_i)-\min_m\widehat T(i,m)]_+/D_w$
and normalize it
over all decision indices $\mathcal I_e$ of the run as
$\nu_t=\Delta_t^{\mathrm{sel}}/\sum_{j\in\mathcal I_e}\Delta_j^{\mathrm{sel}}$.
With failure-credit coefficient $\beta_{\mathrm{miss}}$, the final deadline cost is
\begin{equation}
 g_{D,t}^H=d_t+\beta_{\mathrm{miss}} b_w\omega_t
 +\beta_{\mathrm{miss}} z_e\nu_t,
 \label{eq:method_high_deadline}
\end{equation}
The same coefficient controls both failure-credit terms. Zero denominators
are handled by uniform allocation over the corresponding indices.
Thus, ESP failure credit supplements, rather than replaces, the dense
signal. For queue calibration, let $\mathcal I_{\mathrm{cal}}$ contain
the decision indices in an initial observation window and let
$\theta_q\in(0,1)$ be the desired queue share of the dense cost. Set
\begin{equation}
 \beta_q=\frac{\theta_q}{1-\theta_q}
 \frac{\sum_{t\in\mathcal I_{\mathrm{cal}}}\Delta T_{w,t}/D_w}
 {\sum_{t\in\mathcal I_{\mathrm{cal}}}q_t/(nD_w)}.
 \label{eq:method_queue_calibration}
\end{equation}
The coefficient is then held fixed;
it is zero when no queue waiting is observed.

\paragraph{Reliability.}
The completed workflow reliability deficit is
$\ell_w^R=[\log R_w-\log\mathcal R_w]_+$.
The reliability cost $g_{R,t}^H$ equals $\ell_w^R$ at the last decision
of workflow $w$ and zero at its other decisions. Delayed feedback is
attached to the corresponding transition after completion.

\subsubsection{Value Update}

The three estimators use Double-Q targets~\cite{vanhasselt2016} with a
common next server selected by the online score in
\eqref{eq:method_high_action}. Let this server be $m^*$, let bars denote
target estimators, let $\delta_h$ be discount factors, and let $\iota_t$
indicate the end of the complete scheduling trajectory. Then
\begin{align}
 y_{E,t}&=r_{E,t}^H+\delta_E(1-\iota_t)\bar Q_E(s_{t+1}^H,m^*),\nonumber\\
 y_{D,t}&=g_{D,t}^H+\delta_D(1-\iota_t)[\bar Q_D(s_{t+1}^H,m^*)]_+,\nonumber\\
 y_{R,t}&=g_{R,t}^H+\delta_R(1-\iota_t)\bar Q_R(s_{t+1}^H,m^*).
 \label{eq:method_high_targets}
\end{align}
For replay minibatch $\mathcal B$, each estimator minimizes
\begin{equation}
 \mathcal L_h=\mathbb E_{\mathcal B}\left[
 \operatorname{Huber}\!\left(Q_h(s_t^H,m_i)-y_{h,t}\right)\right],
 \quad h\in\{E,D,R\}.
 \label{eq:method_high_loss}
\end{equation}
Transitions retain the
global chronological order, including switches between workflows,
because their queues and future costs are coupled.

\subsection{Lower-Level Frequency Control}
\label{subsec:low}

\subsubsection{State and Action}

The lower observation $s_t^L$ contains the selected task's workload,
full-frequency execution time, predicted waiting time, remaining
deadline, server power characteristics, and reliability requirement.
It also includes $\chi_t$ and $\chi_t^{\max}$, so that frequency control
can respond to urgency outside the current workflow.
The learner uses SAC~\cite{haarnoja2018}, with a stochastic actor
and two soft critics. The actor produces a squashed-Gaussian
frequency proposal $\rho_i^a\in[\rho_0,1]$, where $\rho_0$ is the
common minimum proposal ratio. The executed ratio is
\begin{equation}
 \rho_i=\max\{\rho_i^a,\rho_{m_i}^{\mathrm{ee}}\}.
 \label{eq:method_frequency_floor}
\end{equation}
The floor excludes frequencies below the energy-efficient operating point; it
does not enforce the time-based frequency in
\eqref{eq:method_time_frequency}.

\subsubsection{Reward and Responsibility Credit}

\paragraph{Energy and local delay.}
For the selected server and the same pre-decision queue, compare the
executed action with full-frequency execution. Define
$\Delta E_i=\widehat E(i,m_i,1)-E_i$ and
$\Delta ET_i=[ET_i-L_i/F_{m_i}]_+$.
If $T_{w,t}^{\mathrm{pre}}$ is the partial workflow response time before
the decision, the local delay shaping cost is
\begin{equation}
 h_t=\frac{\Delta ET_i}{\max\{D_w-T_{w,t}^{\mathrm{pre}},\zeta D_w\}}.
 \label{eq:method_local_delay}
\end{equation}
The energy and delay components are standardized separately using
replay statistics. Write
$\mathcal Z_x(x)=\operatorname{clip}((x-\bar x)/\max\{\sigma_x,\sigma_x^{\min}\},
-b_Z,b_Z)$, where $\bar x$ and $\sigma_x$ are the replay mean and standard
deviation, $\sigma_x^{\min}>0$ is a scale floor, and $b_Z$ bounds the normalized
component.

\paragraph{Deadline responsibility.}
A local timing comparison cannot account for delays passed to later
tasks on the same server. We therefore evaluate completed trajectories
under fixed placement, server execution order, communication delays,
and decision-release times. Restoring execution to full frequency
propagates through both workflow-precedence and server-order edges.
This replay changes execution durations, not the scheduling decisions.

Let $T_v^{(1)}$ be workflow $v$'s response time when all tasks are
restored to full frequency, and $T_v^{(u)}$ its response time when only
workflow $u$ is restored. With normalized lateness
$\mathcal L(T,D)=[(T-D)/D]_+$, define
\begin{align}
 d_v^L&=\min\{1,[\mathcal L(T_v,D_v)-\mathcal L(T_v^{(1)},D_v)]_+\},\nonumber\\
 M_{uv}&=[\mathcal L(T_v,D_v)-\mathcal L(T_v^{(u)},D_v)]_+.
 \label{eq:method_counterfactual_marginal}
\end{align}
The first quantity measures the victim's DVFS-attributable lateness;
the second measures the contribution associated with source workflow
$u$. Allocate each victim's cost back to its sources as
\begin{equation}
 G_u=\sum_{v\in\mathcal W_e}
 \frac{M_{uv}}{\sum_{r\in\mathcal W_e}M_{rv}}\,d_v^L.
 \label{eq:method_source_credit}
\end{equation}
If every marginal for a victim is zero, its cost is allocated in
proportion to the sources' total execution slowdown. With positive
attributable slowdown, $\sum_uG_u=\sum_vd_v^L$, avoiding repeated
charging of the same victim cost. Source credit is distributed among
its frequency actions by
\begin{equation}
 g_{D,t}^L=G_w\frac{\Delta ET_i}
 {\sum_{j\in\mathcal N}\Delta ET_j},\qquad t\in\mathcal I_w,
 \label{eq:method_low_deadline_credit}
\end{equation}
using uniform allocation if the denominator is zero. This conditional
responsibility measure includes queue delays imposed on other
workflows; it is not a reoptimized counterfactual schedule.

\paragraph{Reliability responsibility.}
Let $\widehat{\mathcal R}_{w,t}(f)$ be projected workflow reliability
when the current task executes at frequency $f$, with previous decisions and the
remaining-task estimate held fixed. Define
\begin{align}
 u_{R,t}(f)&=
 \frac{[\log R_w-\log\widehat{\mathcal R}_{w,t}(f)]_+}{-\log R_w},\nonumber\\
 g_{R,t}^L&=[u_{R,t}(f_i)-u_{R,t}(F_{m_i})]_+,
 \label{eq:method_low_reliability}
\end{align}
for $0<R_w<1$. The lower controller is thus charged for the additional
projected deficit caused by its frequency choice, not the deficit
already present at full frequency. Its replay reward is
\begin{equation}
 r_t^L=\mathcal Z_E(\Delta E_i)-\mathcal Z_h(h_t)
 -\mu_D^L g_{D,t}^L-\mu_R^L g_{R,t}^L.
 \label{eq:method_low_reward}
\end{equation}
The local delay term provides shaping, while the two multipliers
regulate responsibility-based constraint costs.

\subsubsection{Actor--Critic Update}

Let $Q_1^L,Q_2^L$ be the critics, $\psi$ the actor parameters,
$\alpha$ the entropy coefficient, and $\delta_L$ the discount factor.
For a sampled next proposal $\rho'{}^a\sim\pi_\psi(\cdot\mid s_{t+1}^L)$,
the soft TD target is
\begin{align}
 y_t^L=r_t^L+\delta_L(1-\iota_t)\big[&\min_{j=1,2}
 \bar Q_j^L(s_{t+1}^L,\rho'{}^a)\nonumber\\
 &-\alpha\log\pi_\psi(\rho'{}^a\mid s_{t+1}^L)\big].
 \label{eq:method_sac_target}
\end{align}
The critics minimize the squared TD error at the stored executed
action $\rho_i$. The actor minimizes
\begin{equation}
 \mathcal L_\pi=\mathbb E\!\left[
 \alpha\log\pi_\psi(\rho^a\mid s^L)
 -\min_{j=1,2}Q_j^L(s^L,\rho^a)\right].
 \label{eq:method_sac_actor}
\end{equation}
Here the policy density refers to the proposal before the operating
floor in \eqref{eq:method_frequency_floor} is applied.

\subsection{Joint Learning and Constraint Adaptation}
\label{subsec:joint}

Both policies interact with the same environment before their updates.
After a complete scheduling trajectory, workflow outcomes determine
the upper failure credits and lower counterfactual credits. These are
attached to chronological transitions and stored in the respective
replay buffers. The upper estimators and the lower actor--critic are
then updated from replay, and their target networks are softly updated.
No additional coordination network is introduced.

Deadline multipliers use proportional--integral constraint feedback,
following the control perspective of Lagrangian
methods~\cite{stooke2020}. Let $\mathcal K_H$ and $\mathcal K_L$ be recent
trajectory windows, and $\vartheta\in(0,1]$ the fraction of the upper
tail used for lower feedback. The two statistics are, respectively,
\begin{equation}
 J_H=\frac{1}{|\mathcal K_H|}\sum_{e\in\mathcal K_H}z_e,
 \qquad
 J_L=\operatorname{CVaR}^{\mathrm{upper}}_{\vartheta}
 \left(\left\{\frac{\sum_{v\in\mathcal W_e}d_v^L}{|\mathcal W_e|}
 \right\}_{e\in\mathcal K_L}\right),
 \label{eq:method_dual_statistics}
\end{equation}
where the empirical upper-tail CVaR is the mean of the largest
$\lceil\vartheta|\mathcal K_L|\rceil$ values. Thus, upper feedback observes
the failure probability underlying ESP, whereas lower feedback
observes the tail of DVFS-attributable lateness, not total system
failure. For level $j\in\{H,L\}$, let $b_j$ be its budget, $\epsilon_j$
the integral step size, $K_j$ the proportional gain, $I_j$ the integral
state, and $\mu_{D,\max}^j$ the multiplier bound. With
$\varepsilon_j=J_j-b_j$, update
\begin{align}
 I_j&\leftarrow\Pi_{[0,\mu_{D,\max}^j]}(I_j+\epsilon_j\varepsilon_j),\nonumber\\
 \mu_D^j&\leftarrow\Pi_{[0,\mu_{D,\max}^j]}(I_j+K_j\varepsilon_j).
 \label{eq:method_pi_dual}
\end{align}
Here $\Pi$ denotes interval projection. The upper training budget
controls a tolerated failure frequency;
strict model-acceptance requirements are assessed separately.

Reliability adaptation uses completed workflow outcomes. Let
$\epsilon_R^H,\epsilon_R^L$ be reliability update step sizes and
$\mu_{R,\max}^H,\mu_{R,\max}^L$ their multiplier bounds. The upper
multiplier accumulates the log-reliability deficit after each completed
workflow, while the lower multiplier compares the recent binary
violation frequency $\bar b_R$ with a small budget $b_R$:
\begin{align}
 \mu_R^H&\leftarrow
 \Pi_{[0,\mu_{R,\max}^H]}(\mu_R^H+\epsilon_R^H\ell_w^R),\nonumber\\
 \mu_R^L&\leftarrow
 \Pi_{[0,\mu_{R,\max}^L]}(\mu_R^L+\epsilon_R^L(\bar b_R-b_R)),
 \label{eq:method_reliability_duals}
\end{align}
where $\bar b_R$ averages $\mathbb{I}(\mathcal R_w<R_w)$ over recent
completed workflows. These updates are intentionally asymmetric: upper reliability
pressure is nondecreasing, whereas lower pressure decreases when its
violation budget is satisfied.

Placement thus learns from end-to-end outcomes, while frequency control
learns from the additional losses attributed to slowdown. The normalized
feedback provides learning surrogates for the physical objective;
service-level feasibility is assessed using the original deadline and
reliability criteria, rather than inferred from the shaped return.

% Preview-only bibliography. In the full paper, use the existing references.bib.
\begin{thebibliography}{4}
\bibitem{wang2016dueling}
Z. Wang, T. Schaul, M. Hessel, H. van Hasselt, M. Lanctot, and
N. de Freitas, ``Dueling network architectures for deep reinforcement
learning,'' in \emph{Proc. ICML}, 2016, pp. 1995--2003.
\bibitem{vanhasselt2016}
H. van Hasselt, A. Guez, and D. Silver, ``Deep reinforcement learning
with double Q-learning,'' in \emph{Proc. AAAI}, vol. 30, no. 1, 2016.
\bibitem{haarnoja2018}
T. Haarnoja, A. Zhou, P. Abbeel, and S. Levine, ``Soft actor-critic:
Off-policy maximum entropy deep reinforcement learning with a stochastic
actor,'' in \emph{Proc. ICML}, 2018, pp. 1861--1870.
\bibitem{stooke2020}
A. Stooke, J. Achiam, and P. Abbeel, ``Responsive safety in reinforcement
learning by PID Lagrangian methods,'' in \emph{Proc. ICML}, 2020,
pp. 9133--9143.
\end{thebibliography}
\end{document}
